% §4 Empirical Validation --- Target 3pp

\section{Empirical Validation}
\label{sec:empirical}

\subsection{Methodology}
\label{sec:emp:method}

A Docker-based regtest harness runs each vault design under a
uniform lifecycle interface; experiments invoke the vault
operations (deposit, trigger, withdraw, recover) without knowledge
of the underlying covenant mechanism and dispatch on per-design
capability flags rather than names. Adapters bind to the
required node variant: Bitcoin Inquisition (\opctv{}, \catcsfs{}),
Merkleize \opccv{}~\cite{BIP443}, jamesob's \opvault{} branch, and
Elements under Simplex (Simplicity, Appendix~\ref{app:simplicity}).
The primary metric is virtual size (vsize): structurally
deterministic, identical on regtest and mainnet. Fees are
projections $\text{fee} = \text{vsize} \times r$ since regtest has
no fee market. Temporal recovery races are not captured (blocks mine
on demand). Fifteen experiments span lifecycle costs, batching,
revault chains, address reuse, class \classfee{} (fee-channel
pinning), class \classscale{} (per-UTXO recovery-cost scaling), and
the 11-element threat taxonomy; measurements are reproducible via the
anonymised artifact (DOI withheld for double-blind review).

\subsection{Lifecycle Transaction Sizing}
\label{sec:emp:sizes}

Table~\ref{tab:vsizes} reports the measured vsize for every
transaction in the deposit--trigger--withdraw lifecycle. Each value
is structurally deterministic (range~=~0 across repeated runs). The
rightmost column tags the axis value whose consequences each row
measures, per Section~\ref{sec:bg:axes}.

\begin{table}[t]
\centering
\caption{Lifecycle vsize measurements (\vB{}) for the four
Bitcoin covenant proposals. Lifecycle total~=~deposit +
trigger + withdraw. Simplicity on Elements is presented in
Appendix~\ref{app:simplicity}. The \emph{axis-value} column gives
the primary axis (Table~\ref{tab:design_space}) whose position the
row's size measures as a consequence.}
\label{tab:vsizes}
\begin{tabular}{@{}lrrrrcl@{}}
\toprule
Design      & deposit & trigger & withdraw & recover & \textbf{total} & axis-value                                          \\
\midrule
\opctv{}    & 122     &  94     & 152      & 133     & \textbf{368}   & \axfee{}=anchor, \axwdbind{}=digest                 \\
\opccv{}    & 300     & 154     & 111      & 122     & \textbf{565}   & \axrevault{}=yes, \axfee{}=dynamic                  \\
\opvault{}  & 154     & 292     & 121      & 246     & \textbf{567}   & \axrevault{}=yes, \axfee{}=dynamic (wallet)         \\
\catcsfs{}  & 122     & 221     & 210      & 125     & \textbf{553}   & \axwdbind{}=dual-bind, \axfee{}=dynamic (ACP)       \\
\bottomrule
\end{tabular}
\end{table}

\opctv{} has the cheapest lifecycle (368\vB{}); its trigger is a bare
template-hash check with no signature (94\vB{}, smallest of any
design). \opccv{} and \opvault{} total nearly identically (565 vs.\
567\vB{}) with opposite distributions: \opccv{} has a cheap trigger
(154\vB{}) and expensive deposit (300\vB{}, two contract outputs);
\opvault{} has an expensive trigger (292\vB{}, fee-wallet input) and
moderate deposit. \catcsfs{}'s dual \opcsfs{}+CHECKSIG pattern falls
between (553\vB{}), costing trigger 221\vB{} and withdraw 210\vB{}.

\paragraph{\opvault{} vsize correction.}
Prior informal estimates of \opvault{} \texttt{start\_withdrawal}
(${\sim}200$\vB{}) and recovery (${\sim}170$\vB{}) are both under by
$\approx 45\%$ (measured: 292\vB{}, 246\vB{}). BIP-345's
value-preservation rule forces a P2TR fee-wallet input
($80$--$90$\vB{}) on every non-deposit transaction, yielding a 36\%
lifecycle cost premium over \opccv{}---an economic factor in the BIP-345
withdrawal~\cite{OB25}.

\subsection{Class Measurements: \classfee{} and \classscale{}}
\label{sec:emp:attacks}
\label{sec:empirical:classes}

We reproduce the two primary economic attack classes on regtest. For
class \classfee{} (fee-channel pinning, TM1,
Section~\ref{sec:imposs:class_fee}), we construct the 25-chain
descendant attack against a \opctv{} unvault's 3{,}300-sat anchor and
confirm $C_{\text{pin}} \approx 14{,}300$~sats---$0.03\%$ of a
$0.5$~BTC vault, rational at any fee regime. For class \classscale{}
(per-UTXO recovery-cost scaling, TM4,
Section~\ref{sec:imposs:class_scale}), substituting the measured
$\text{vsize}_{\text{rec}}$ of Table~\ref{tab:vsizes} into
Eq.~\ref{eq:r-dagger} yields unbatched thresholds
$r^\dagger_{\mathrm{OP\_CCV}} \approx 66$\satvB{} and
$r^\dagger_{\opvault{}} \approx 59$\satvB{} at $V = 0.5$~BTC. Above
$r^\dagger$ the structural cost tail saturates within a single block;
below it the defender can absorb the per-UTXO liability before
dust-economics kick in. \opctv{} and \catcsfs{} are immune to
\classscale{} by \axrevault{} = \texttt{no} (atomic withdrawal only).

\subsection{Batched Recovery and the Ordering Flip}
\label{sec:empirical:batching}

This is the paper's empirically novel claim: under class \classscale{},
the safety ranking between \opccv{} and \opvault{} inverts based on
whether the watchtower batches recovery. The multi-round attacker
scenario of prior framings is one threat-model instantiation
among several (Appendix~\ref{app:rationality}); the ordering-flip
result holds independently of attacker behaviour and is a statement
about \emph{defender implementation choice}.

BIP-345 §``Batching'' and BIP-443 default aggregation permit the
defender to recover $N$ triggered UTXOs in a single transaction. We
swept $N \in \{2, 5, 10, 25, 50\}$ on our regtest harness for each
batching-capable design; linear regression yields
$\text{vsize}_{\text{rec}}(D, N) = \alpha(D) + \beta(D) \cdot N$:
\begin{itemize}
  \item \opccv{}: $\alpha_{\mathrm{OP\_CCV}} = \overheadCCV$~\vB{},
    $\beta_{\mathrm{OP\_CCV}} = \perinputCCV$~\vB{}.
  \item \opvault{}: $\alpha_{\opvault{}} = \overheadOPV$~\vB{},
    $\beta_{\opvault{}} = \perinputOPV$~\vB{}.
\end{itemize}

Substituting $\beta(D)$ for $\text{vsize}_{\text{rec}}(D)$ in
Eq.~\ref{eq:r-dagger}, the batched thresholds at $V = 0.5$~BTC rise
to 119\satvB{} (\opccv{}) and 159\satvB{} (\opvault{}), factors of
$1.8\times$ and $2.7\times$ over unbatched. The ordering flips: \opccv{}
is safer unbatched ($66 > 59$) but \opvault{} is safer batched
($159 > 119$). \opvault{}'s larger unbatched recovery transaction
amortises disproportionately; design choice under \classscale{} thus
depends on the defender's batching strategy, and an ordering stated
without reference to $b$ is incomplete.

\paragraph{Asymmetric benefit.}
BIP-345 authorised recovery permits batching across vaults with
distinct \texttt{<recovery-sPK-hash>} values; unauthorised recovery
is capped at two outputs and thus restricted to
same-\texttt{recovery-sPK} batches. \opccv{}'s default aggregation has no
taptree-equivalence requirement and is strictly more permissive. The
attacker's partial-unvault-with-revault sequence is inherently
serial---each revault consumes the previous output---so the
attacker does not benefit from batching, and the measured savings
accrue only to the defender.

\subsection{Simplicity on Elements}
\label{sec:simplicity-perspective}

Our Simplicity vault on Elements supports per-output binding on
trigger and recovery, atomic or bounded partial withdrawal, cold-key
recovery with an optional keyless post-timeout sweep, and MAX\_FEE-capped
in-value fee deduction. In the atomic configuration with
cold-key-only recovery, it occupies the dominant corner of
Section~\ref{sec:imposs:design_space} (strict-best on \axauth{},
\axrevault{}, \axwdbind{}, \axrecbind{}, and \axintro{}
simultaneously). Empirical measurements (lifecycle 865\vB{} for the
atomic configuration; threat coverage; Coq-verified semantics) and
the bounded-partial and keyless-post-timeout constructions are
reported in Appendix~\ref{app:simplicity}. Its role in the argument
is to demonstrate that the dominant corner is occupiable on some
substrate, even while no Bitcoin proposal currently occupies it.

\providecommand{\todo}[2][]{\textbf{[TODO: #2]}}
